\subsection{Fixed-Valency Comparisons Across Admissible Groups}
\label{subsec:fixed-valency-comparisons}

A potential objection to the comparison of $R_p$ values across different groups is
that different Cayley graphs may have different valencies (numbers of generators),
making a direct comparison of $\lambda_\rho$ values unfair.
We address this by comparing groups at fixed valency $|S|=6$.

For valency $|S|=6$, the admissible groups and their spinorial sector eigenvalues are:

\begin{center}
  \begin{tabular}{lcccc}
    \hline
    Group & $|S|$ & $\lambda_{\rho_s}$ & $\lambda_{\rho_{\mathrm{ab}}}^{\min}$
    & $R_p$ \\
    \hline
    $Q_8$ & 6 & 6 & 8 & $\sqrt{4/3}\approx 1.155$ \\
    $\mathrm{Dic}_3$ & 6 & 6 & 8 & $\approx 1.155$ \\
    $2T$ (excluded) & 6 & --- & --- & not admissible \\
    \hline
  \end{tabular}
\end{center}

At valency 6, $Q_8$ and $\mathrm{Dic}_3$ are the minimal admissible structures,
and they share the same admissibility ratio.
The exclusion of $2T$ at valency 6 is confirmed by the Gram constraint failure
documented in D.13.

For higher valencies, larger binary polyhedral groups become accessible.
The general trend is that $R_p$ increases with group complexity within the binary
family, before converging to 1 as the continuum limit is approached.
This confirms that the spectral admissibility filter is a genuinely resolution-dependent
phenomenon, with the binary polyhedral groups forming an ordered hierarchy of
increasingly refined finite approximations to the $SU(2)$ attractor.
